\chapter{Implementation snapshot (aligned with code)}

\section{Entry point and trainer}
\texttt{train\_battler.py} parses CLI flags, loads \texttt{TrainingConfig} from \texttt{get\_config(preset)}, sets the MLflow experiment name \texttt{Pokemon\_RL\_Battler}, and runs \texttt{PokemonTrainer.train()} (\texttt{src/training/trainer.py}). The trainer registers the env, builds RLlib PPO config (\texttt{rllib\_config\_builder.py}), manages checkpoints, curriculum updates via callbacks, and logs metrics.

\section{Parallelism and Showdown}
\texttt{EnvironmentConfig} defaults: \texttt{num\_workers=12}, \texttt{num\_envs\_per\_worker=4} (48 envs), \texttt{num\_servers=8}, \texttt{start\_port=8000}. Ports are assigned deterministically to workers. Eight Node servers should be running before training (\texttt{./scripts/spin\_up\_multiple\_showdown.sh}).

\section{Battle format}
Default \texttt{battle\_format} is \texttt{gen5randombattle} (see \texttt{TM\_optimal\_config.py}). This can differ from older docs that mention Gen~8 random doubles/singles; always read the config field for the experiment you run.

\section{Curriculum stages (defaults)}
Table~\ref{tab:curr} lists \texttt{CurriculumConfig} defaults in \texttt{TM\_optimal\_config.py}. Promotion uses a deque of win/loss outcomes of length \texttt{rolling\_window\_episodes=200}, requires at least \texttt{min\_episodes\_before\_promotion=100} in-stage episodes, and compares rolling win rate to \texttt{promote\_at\_win\_rate}.

\begin{table}[h]
\centering
\small
\begin{tabularx}{\textwidth}{l X c c}
\toprule
\textbf{Stage} & \textbf{Opponent mix} & \textbf{Promote at} & \textbf{Min samples} \\
\midrule
\texttt{random\_warmup} & \texttt{random\_no\_switch: 1.0} & 0.70 & 100 \\
\texttt{self\_play} & \texttt{self: 0.8}, \texttt{random\_no\_switch: 0.2} & 0.85 & 200 \\
\texttt{mixed} & \texttt{heuristic: 0.5}, \texttt{self: 0.3}, \texttt{random\_no\_switch: 0.2} & 1.01 (terminal) & 200 \\
\bottomrule
\end{tabularx}
\caption{Default curriculum stages (subject to change in config).}
\label{tab:curr}
\end{table}

\section{Opponent keys}
\begin{itemize}
  \item \texttt{random} --- poke-env \texttt{RandomPlayer} (uniform legal orders, including switches).
  \item \texttt{random\_no\_switch} --- custom \texttt{RandomNoSwitchPlayer}: random among legal \emph{move} orders when available; switches only when forced by faint or no moves (\texttt{src/envs/random\_no\_switch\_player.py}).
  \item \texttt{heuristic} --- \texttt{SimpleHeuristicsPlayer}.
  \item \texttt{self} --- \texttt{SelfPlayPlayer} with exported weights path from config.
\end{itemize}

In observations, \texttt{random\_no\_switch} is \textbf{canonicalized} to the same global opponent-type signal as \texttt{random} in \texttt{embedding.py} to limit observation shift when mixing baselines.

\section{Reward defaults}
Base \texttt{RewardConfig}: victory/defeat $\pm 10$; HP weight $2.0$; faint reward/penalty $\pm 5$; step penalty $-0.005$; matchup and action-quality weights $0.2$ and $0.3$. Per-stage copies in the curriculum use the same numeric defaults in the current file.

\section{Model configuration (defaults)}
\texttt{ModelConfig}: \texttt{num\_tokens=13}, \texttt{token\_dim=164}, embedding dim 32, transformer \texttt{hidden\_dim=512}, \texttt{num\_heads=8}, \texttt{num\_transformer\_layers=1}, \texttt{use\_position\_embeddings=True}, \texttt{use\_role\_embeddings=True}, \texttt{use\_lstm=True}, \texttt{lstm\_hidden=512}. Vocabulary sizes come from \texttt{data/vocab/gen8\_embedding\_vocab.json} via \texttt{vocab\_sizes()}. Presets override some fields (e.g.\ \texttt{quick} and \texttt{memory\_safe} set \texttt{use\_lstm=False}).

\section{PPO defaults (dataclass)}
Illustrative fields: \texttt{lr}$=3\times10^{-4}$, \texttt{gamma}$=0.95$, \texttt{lambda\_}$=0.95$, \texttt{clip\_param}$=0.2$, \texttt{entropy\_coeff}$=0.05$, \texttt{vf\_loss\_coeff}$=0.5$, \texttt{grad\_clip}$=5.0$, batch \texttt{train\_batch\_size}$=6144$, \texttt{sgd\_minibatch\_size}$=256$, \texttt{num\_sgd\_iter}$=10$.

\section{Checkpoints and validation}
Default checkpoint directory \texttt{checkpoints/}, frequency \texttt{checkpoint\_freq=250\,000}, retaining 5 checkpoints. Scheduled validation (\texttt{ValidationScheduleConfig}) can run protocols such as smoke, fixed paired teams, and mirror matches against manifests under \texttt{data/validation/}.

\section{LSTM and RLlib connectors}
The transformer module supports an LSTM head for cross-turn memory when \texttt{use\_lstm=True}. The code documents connector constraints: enabling LSTM requires the rollout pipeline to carry recurrent state correctly. Presets that disable LSTM sidestep this for stability during development.

\section{Historical note (\texttt{temp.md} week-7 text)}
\texttt{temp.md} describes an older three-stage narrative (easy/medium/hard with 85\%/75\% thresholds and $\pm80$/$\pm100$ terminal rewards). That text is useful historically but \textbf{does not match} the curriculum table above; comparisons across time must name the config commit or file version.
